%=====================================================================
\chapter{The Artificial Intelligence Project}\label{ch:project}
%=====================================================================
\locator{6}{Part VI --- Learning, Language and Responsible AI}

\begin{outcomes}
By the end of this chapter you will be able to:
\begin{tight}
  \item \textbf{Choose} and \textbf{scope} a problem that the methods of this
        course can actually solve.
  \item \textbf{Write} a proposal stating the PEAS specification, the method and
        the measure of success.
  \item \textbf{Implement} a working solution with a baseline to compare it
        against.
  \item \textbf{Report} the result honestly, including what did not work.
  \item \textbf{Audit} your own system as \chref{responsible} requires.
\end{tight}
\end{outcomes}

\begin{prereq}
The whole course. The project is set in Week~11 and presented in Week~16
(Appendix~G); the proposal is due the week it is set.
\end{prereq}

\begin{keyterms}
problem scoping $\bullet$ PEAS specification $\bullet$ proposal $\bullet$ baseline
$\bullet$ evaluation measure $\bullet$ ablation $\bullet$ reproducibility $\bullet$
project report $\bullet$ demonstration $\bullet$ audit
\end{keyterms}

\section{The Brief}

Build a working AI system for a problem you can describe in one sentence, using one
or more of the methods from this course. Demonstrate it, report honestly on what it
does and does not do, and audit it.

\begin{examplebox}[What is being assessed]
\begin{tabular}{@{}L{4.2cm}L{1.8cm}L{7.8cm}@{}}
\toprule
\textbf{Component} & \textbf{Weight} & \textbf{What it means} \\
\midrule
Proposal & 10\% & PEAS, method, measure, baseline --- before any code \\
Implementation & 30\% & It runs, it is reproducible, somebody else can read it \\
Evaluation & 25\% & Compared against a baseline, on a measure you defined in
advance \\
Report & 25\% & Including what failed and why \\
Demonstration & 10\% & Ten minutes, live, with one hard question answered \\
\bottomrule
\end{tabular}

\smallskip
Note the weights: the code is under a third. A working system with no baseline and
no honest evaluation scores worse than a modest one with both.
\end{examplebox}

\section{Choosing a Problem}

The commonest failure is not technical. It is choosing a problem whose difficulty
lies somewhere this course did not go --- almost always in acquiring data or in
integrating with something.

\begin{conceptbox}
Three tests for a candidate problem, in order.

\emph{Can you state the goal in one sentence, without the word ``and''?} ``Schedule
a sprint respecting precedence and capacity'' passes. ``Predict delays and route
tickets and suggest owners'' is three projects.

\emph{Do the rules exist, and can you write them down in an afternoon?} This course
builds agents that are \emph{told} the rules. If the rules would have to be learned
from data, you have chosen a \emph{Machine Learning} project and you do not yet have
the machinery.

\emph{Is there a baseline you would be embarrassed to lose to?} First-come
first-served, a greedy heuristic, or the current manual process. If you cannot name
one, you will not be able to show that your system helped.
\end{conceptbox}

\smallskip
\noindent\begin{longtable}{@{}L{5.0cm}L{2.5cm}L{7.0cm}@{}}
\toprule
\rowcolor{primary!15}
\textbf{Problem} & \textbf{Method} & \textbf{Baseline to beat} \\
\midrule
\endhead
Timetable a set of lab sessions against rooms and staff & \chref{csp} &
Greedy assignment in enrolment order \\
Plan a deployment from operator descriptions & \chref{planning} & The runbook's
fixed order \\
Route service-desk tickets to a queue & \chref{expert} & The current manual
triage \\
Diagnose the likely cause of a recurring failure & \chref{bayesnets} & Always blame
the most common cause \\
Choose which of a backlog to take into a sprint & \chref{informed} & Highest
priority first \\
Rate risk for a portfolio of projects & \chref{fuzzy} & A threshold rule \\
Play a two-player scheduling game against a tutor & \chref{adversarial} & Random
legal move \\
Answer questions about a domain from a knowledge base & \chref{fol} &
Keyword search \\
\bottomrule
\caption{Eight projects that fit the course, each with a baseline you can implement
in an hour. Your own problem is welcome and must pass the three tests
above.}\label{tab:projects}
\end{longtable}

\section{The Proposal}

One page, due the week the project is set. Five headings, and the fifth is the one
students leave out.

\begin{enumerate}[leftmargin=2.2em, itemsep=3pt, topsep=4pt]
  \item \textbf{The problem}, in one sentence.
  \item \textbf{PEAS} (\chref{agents}): performance measure, environment,
        actuators, sensors. The performance measure must describe the state of the
        environment, not the activity of your program.
  \item \textbf{Method}, and why. Name the chapter. If two methods could work, say
        which you chose and what you gave up.
  \item \textbf{Baseline}, and how you will implement it.
  \item \textbf{What would count as failure.} Write down, in advance, the result
        that would make you say the approach did not work. This is the heading that
        turns an evaluation into evidence rather than a demonstration.
\end{enumerate}

\begin{alertbox}[Decide the measure before you build]
If the measure is chosen afterwards, it will be chosen --- unconsciously and
honestly --- to be one the system does well on. That is not dishonesty; it is how
people work, and the only defence is to write the measure down first and then not
change it.

If you do change it, say so in the report and say why. A report that says ``we
originally measured X, found it was the wrong question for this reason, and switched
to Y'' is \emph{stronger} than one that never mentions it, because it shows you
understood what you were measuring.
\end{alertbox}

\section{Building It}

\begin{tight}
  \item \textbf{The baseline first.} It takes an hour and it is the only thing that
        makes your result interpretable. Build it before the interesting part, or
        you will not build it at all.
  \item \textbf{Seed everything.} Any randomness gets an explicit seed, as every lab
        in this book does. A result nobody can reproduce is not a result.
  \item \textbf{Keep the domain separate from the algorithm.} The constraints, the
        operators, the rules go in data. This is \chref{introduction}'s third test,
        it is what makes the system correctable, and it is visibly assessed.
  \item \textbf{Log a trace.} Whatever your method, record why it did what it did.
        You will need it for the report, for the demonstration's hard question, and
        for the audit.
  \item \textbf{Measure at the size you will demonstrate at.} \chref{csp},
        \chref{planning} and \chref{bayesnets} each contained a technique that
        \emph{lost} on a small instance and won on a large one. Do not tune on a toy.
\end{tight}

\section{Reporting It}

Eight to twelve pages. The structure below is not a suggestion; reports that depart
from it lose marks for the sections they omit.

\smallskip
\noindent\begin{longtable}{@{}L{3.4cm}L{11.2cm}@{}}
\toprule
\rowcolor{primary!15}
\textbf{Section} & \textbf{What goes in it} \\
\midrule
\endhead
Problem and PEAS & One page. The same specification as the proposal, corrected if
it changed \\
Method & Which chapter's machinery, and why this one. State the alternative you
rejected \\
Implementation & Representation first, algorithm second. A figure of the data
structure is worth a page of prose \\
Evaluation & The baseline, your system, on the measure you fixed in advance. A
table and one figure \\
What did not work & Required. Every project has some, and its absence is read as
not having looked \\
Audit & \chref{responsible}: outcomes by group where that applies, and the
situations your system does not cover \\
Limits & Where it breaks, and the size at which it stops being usable. Be specific
and quantitative \\
\bottomrule
\caption{The report structure. The last three sections are what separate a project
from a demonstration.}\label{tab:report}
\end{longtable}

\begin{conceptbox}
On the \emph{what did not work} section, since it is the one people resist.

Every project in this course will have a component that failed, an approach
abandoned, or a measurement that came out the wrong way. Reporting it costs nothing
--- it is not marked as failure --- and omitting it costs a great deal, because a
report in which everything worked first time is not believed by anybody who has
built anything.

This book has tried to model the habit. \chref{csp}'s ordering heuristic made things
worse; \chref{planning}'s heuristic bought under one per cent; \chref{bayesnets}'
variable elimination lost on the small network; \chref{responsible}'s ablation was
fooled by a rule whose removal abolished the outcome. Each of those is a real
measurement that contradicted the tidy expectation, and each taught more than the
tidy version would have.
\end{conceptbox}

\section{The Demonstration}

Ten minutes in Week~16. Run the system live on an input the audience chooses within
the constraints you state. Show the trace. Then answer one hard question, which will
be some version of: \emph{how do you know it is right?}

Prepare for that question specifically. The good answer names the baseline, the
measure and the margin, and says where the system stops working.

\begin{worked}[A proposal that would score well]
\textbf{Problem.} Assign this department's laboratory sessions to rooms and
demonstrators without clashes.

\smallskip
\textbf{PEAS.} \emph{Performance measure:} number of sessions placed, zero hard
constraint violations, and the spread of load across demonstrators measured by the
difference between the busiest and least busy. Note this describes the timetable,
not the program. \emph{Environment:} the session list, room capacities, staff
availability, the week. Fully observable, deterministic, static for the duration of
a solve, discrete, single-agent. \emph{Actuators:} assign a session to a
(room, slot, demonstrator) triple; unassign. \emph{Sensors:} the full session list
and all constraints.

\smallskip
\textbf{Method.} Constraint satisfaction (\chref{csp}). Variables are sessions,
domains are legal triples, constraints are room capacity, no double-booking of a
room or a person, and precedence where a lab must follow its lecture. Backtracking
with minimum remaining values and forward checking --- and \chref{csp}'s measured
result says the propagation is what will make MRV worth having.

\emph{Alternative rejected:} local search (\chref{local}) would produce a good
timetable faster, and could not \emph{prove} the absence of a clash. A timetable
with a clash is worthless however good its load balance, so the guarantee matters
more than the speed here.

\smallskip
\textbf{Baseline.} Greedy assignment in enrolment order, taking the first legal
triple. About thirty lines.

\smallskip
\textbf{What would count as failure.} If the greedy baseline places as many sessions
with a load spread within two of mine, constraint satisfaction has not earned its
complexity on this instance and I will say so. If the solver cannot place the full
set within sixty seconds, I will report the largest instance it does handle.

\smallskip
\textbf{Why this proposal is good.} The goal is one sentence. The rules exist and
are already written down in the department handbook. There is an embarrassing
baseline. The performance measure describes the world. And the failure condition is
stated in advance, in numbers, so the evaluation cannot be quietly reshaped to
flatter the result.
\end{worked}

\begin{pitfall}
\begin{tight}
  \item \textbf{Choosing a problem whose difficulty is the data.} If the rules must
        be learned, it is a \emph{Machine Learning} project.
  \item \textbf{No baseline.} Then your result has no units. Build it first.
  \item \textbf{Choosing the measure afterwards.} It will be a measure your system
        does well on, and you will not notice.
  \item \textbf{Knowledge inside the algorithm.} Constraints and rules belong in
        data. This is assessed, and it is what makes your system correctable.
  \item \textbf{No trace.} You will need it for the report, for the hard question,
        and for the audit. Retrofitting one is much harder than logging as you go.
  \item \textbf{Tuning on a toy instance.} Three chapters of this book contain a
        technique that loses on a small instance and wins on a large one.
  \item \textbf{Omitting ``what did not work''.} It is required, it costs no marks,
        and its absence is conspicuous.
\end{tight}
\end{pitfall}

%---------------------------------------------------------------------
\begin{lab}[The Project Template]
The scaffold to start from. It runs as given, with a trivial problem in place, so
that the structure is working before any of your own code exists.
\begin{lstlisting}[language=Python]
"""Chapter 22 -- the project template.

Replace the marked sections. Everything else is the shape every
project in this course should have: a PEAS specification you can
print, a baseline, a measure fixed in advance, a trace, an audit and
a reproducible run.
"""

import json
import os
import random

SEED = 0                      # every random choice, reproducible

# ===================================================================
# 1. PEAS -- fill this in, and print it in the report
# ===================================================================
PEAS = {
    "performance": "sessions placed; zero clashes; load spread",
    "environment": "sessions, rooms, staff, slots; static, discrete",
    "actuators": "assign(session, room, slot, staff); unassign",
    "sensors": "the full session list and all constraints",
}

# what would count as FAILURE -- written down BEFORE building
FAILURE_CONDITION = (
    "the baseline places as many sessions with a load spread within "
    "2 of ours"
)

# ===================================================================
# 2. The domain, as DATA -- never inside the algorithm
# ===================================================================
SESSIONS = ["ai-lab-1", "ai-lab-2", "ds-lab-1", "ml-lab-1"]
SLOTS = ["mon-am", "mon-pm", "tue-am"]
STAFF = ["ayesha", "bilal"]


def legal(assignment, session, slot, who):
    """The constraints. Reading this should tell a domain expert
    exactly what the system believes."""
    for (s, (sl, w)) in assignment.items():
        if sl == slot and w == who:
            return False              # one person, one slot
    return True


# ===================================================================
# 3. The BASELINE -- build this first, always
# ===================================================================
def baseline(sessions):
    """Greedy: first legal option, in the order given."""
    assignment = {}
    for s in sessions:
        for slot in SLOTS:
            for who in STAFF:
                if legal(assignment, s, slot, who):
                    assignment[s] = (slot, who)
                    break
            if s in assignment:
                break
    return assignment


# ===================================================================
# 4. YOUR SYSTEM -- replace this with the real method
# ===================================================================
def solve(sessions, trace):
    """Backtracking with minimum remaining values would go here.

    The template uses plain backtracking so the scaffold runs; the
    point is that it LOGS, and that the log is what the report,
    the demonstration and the audit are all built from.
    """
    assignment = {}

    def bt(remaining):
        if not remaining:
            return True
        # MRV would order 'remaining' here
        s = remaining[0]
        for slot in SLOTS:
            for who in STAFF:
                if legal(assignment, s, slot, who):
                    assignment[s] = (slot, who)
                    trace.append("assign %s -> %s, %s" % (s, slot, who))
                    if bt(remaining[1:]):
                        return True
                    trace.append("undo %s" % s)
                    del assignment[s]
        return False

    bt(list(sessions))
    return assignment


# ===================================================================
# 5. The MEASURE -- fixed in advance, and the same for both
# ===================================================================
def measure(assignment):
    placed = len(assignment)
    load = {w: 0 for w in STAFF}
    for (_, who) in assignment.values():
        load[who] += 1
    spread = max(load.values()) - min(load.values()) if load else 0
    return {"placed": placed, "spread": spread, "load": load}


# ===================================================================
# 6. The AUDIT -- Chapter 21
# ===================================================================
def audit(assignment):
    """Outcomes by group, and the situations not covered."""
    by_course = {}
    for s, (slot, who) in assignment.items():
        course = s.split("-")[0]
        by_course.setdefault(course, []).append(slot)
    unplaced = [s for s in SESSIONS if s not in assignment]
    return {"by_course": by_course, "unplaced": unplaced}


if __name__ == "__main__":
    random.seed(SEED)
    print("=" * 58)
    print("PEAS")
    print("=" * 58)
    for k, v in PEAS.items():
        print("   %-13s %s" % (k + ":", v))
    print()
    print("   fails if:    %s" % FAILURE_CONDITION)

    print()
    print("=" * 58)
    print("BASELINE vs SYSTEM")
    print("=" * 58)
    base = baseline(SESSIONS)
    trace = []
    ours = solve(SESSIONS, trace)

    mb, mo = measure(base), measure(ours)
    print("   measure      baseline   system")
    print("   " + "-" * 34)
    for key in ("placed", "spread"):
        print("   %-12s %8s %8s" % (key, mb[key], mo[key]))
    print("   load         %8s %8s" % (mb["load"], mo["load"]))

    # state the verdict against the CONDITION FIXED IN ADVANCE
    print()
    beat = (mo["placed"] > mb["placed"]
            or (mo["placed"] == mb["placed"]
                and mo["spread"] < mb["spread"] - 2))
    print("   verdict: %s"
          % ("the system beat the baseline"
             if beat else
             "the system did NOT beat the baseline on the measure "
             "fixed in advance -- report this"))

    print()
    print("=" * 58)
    print("TRACE (first 8 steps)")
    print("=" * 58)
    for line in trace[:8]:
        print("   " + line)
    print("   ... %d steps total" % len(trace))

    print()
    print("=" * 58)
    print("AUDIT")
    print("=" * 58)
    a = audit(ours)
    for course, slots in sorted(a["by_course"].items()):
        print("   %-6s %d session(s): %s"
              % (course, len(slots), ", ".join(slots)))
    print("   unplaced:", a["unplaced"] or "none")

    # reproducibility: the artefact the report cites
    out = {"seed": SEED, "peas": PEAS,
           "failure_condition": FAILURE_CONDITION,
           "baseline": mb, "system": mo, "audit": a,
           "trace_length": len(trace)}
    os.makedirs("project_output", exist_ok=True)
    with open(os.path.join("project_output", "metrics.json"), "w") as f:
        json.dump(out, f, indent=2, default=str)
    print()
    print("wrote project_output/metrics.json -- cite this in the")
    print("report, and commit it. A result nobody can reproduce is")
    print("not a result.")

    # the scaffold's own checks
    assert mo["placed"] == len(SESSIONS), "the system should place all"
    assert os.path.exists(os.path.join("project_output",
                                       "metrics.json"))
\end{lstlisting}
Run it before writing any of your own code. If the scaffold works, you have a
baseline, a measure, a trace, an audit and a reproducible artefact on day one --- and
everything after that is the interesting part.

\smallskip
Notice what the template prints as it stands: \emph{the system did NOT beat the
baseline on the measure fixed in advance}. That is correct --- on four sessions the
greedy baseline is already optimal --- and it is the behaviour to keep. A scaffold
that congratulated itself before anybody had written a method would be training you
in exactly the habit this chapter is trying to prevent.
\end{lab}

%---------------------------------------------------------------------
\begin{checkpoint}
\begin{tightnum}
  \item A student proposes ``an AI system that predicts which tickets will breach
        their service level agreement''. Apply the three tests of this chapter and
        say what you would advise.
  \item Why must the measure and the failure condition be written down before the
        system is built? Give the failure mode that results from not doing so.
  \item Your system beats the baseline by a small margin. Name three things the
        report must contain before that claim means anything.
\end{tightnum}
\end{checkpoint}

%---------------------------------------------------------------------
\begin{chaptersummary}
\begin{tight}
  \item Choose a problem whose goal is one sentence, whose rules already exist, and
        which has a baseline you would be embarrassed to lose to. If the rules must
        be learned from data, it is a \emph{Machine Learning} project.
  \item The proposal is five headings, and the fifth --- \emph{what would count as
        failure} --- is what turns an evaluation into evidence.
  \item Build the baseline first, seed everything, keep the domain in data, log a
        trace, and measure at the size you will demonstrate at.
  \item The report has seven sections. \emph{What did not work}, the \emph{audit}
        and the \emph{limits} are what separate a project from a demonstration.
  \item Decide the measure before building. Chosen afterwards, it will be one the
        system does well on, and honestly so.
  \item Prepare for one question at the demonstration: \emph{how do you know it is
        right?} The good answer names the baseline, the measure, the margin and
        where the system stops working.
\end{tight}
\end{chaptersummary}

\begin{reviewq}
  \item \textbf{(MCQ)} The first thing to build in a project is:
        (a)~the main algorithm (b)~the baseline (c)~the user interface
        (d)~the report.
  \item \textbf{(MCQ)} A performance measure should describe:
        (a)~how much the program does (b)~the state of the environment
        (c)~the running time (d)~the number of rules.
  \item \textbf{(MCQ)} ``What would count as failure'' is written:
        (a)~in the report (b)~in the proposal, before building
        (c)~after evaluation (d)~only if the project fails.
  \item \textbf{(MCQ)} A project whose rules must be learned from examples is:
        (a)~ideal for this course (b)~a Machine Learning project
        (c)~a constraint problem (d)~impossible.
  \item \textbf{(Short)} Give a one-sentence problem statement for a project of
        your own, and name the chapter whose method you would use and the baseline
        you would beat.
  \item \textbf{(Short)} Explain why keeping the domain in data rather than in the
        algorithm is assessed, referring to \chref{introduction}'s checklist.
  \item \textbf{(Short)} Your evaluation shows the baseline winning. Describe how to
        report that so it is a good report rather than a failed project.
  \item \textbf{(Applied)} Write the full five-heading proposal for one of the
        problems in Table~\ref{tab:projects}, including the PEAS specification, the
        alternative method you reject with a reason, the baseline, and a failure
        condition stated in numbers.
\end{reviewq}
